\documentclass[11pt]{article}
\newcommand{\LabNumber}{8}
\newcommand{\LabTitle}{Trees and Tree Traversals}
\newcommand{\LabTopic}{Binary tree, preorder/inorder/postorder, depth, height}
\input{common.tex}

\begin{document}
\labheader

\section*{Objectives}
\begin{itemize}[leftmargin=*,nosep]
  \item Implement a binary tree using a linked-node structure with \texttt{left} and \texttt{right} children.
  \item Implement preorder, inorder, and postorder traversals recursively.
  \item Compute the depth of a node and the height of a tree recursively.
  \item Count nodes, count leaves, and evaluate an expression tree.
\end{itemize}

\begin{summary}[Quick Reference]
\textbf{Node class.}
\begin{lstlisting}
class Node<E> {
    E element;
    Node<E> left, right, parent;
    Node(E e, Node<E> p) { element = e; parent = p; }
}
\end{lstlisting}

\textbf{Traversals.}
\begin{lstlisting}
void preorder(Node<E> v)  { visit(v); preorder(v.left);  preorder(v.right); }
void inorder(Node<E> v)   { inorder(v.left);  visit(v); inorder(v.right);  }
void postorder(Node<E> v) { postorder(v.left); postorder(v.right); visit(v); }
\end{lstlisting}

\textbf{Height and depth.}
\begin{lstlisting}
int height(Node<E> v) {
    if (v == null) return -1;
    return 1 + Math.max(height(v.left), height(v.right));
}
int depth(Node<E> v) { return v.parent == null ? 0 : 1 + depth(v.parent); }
\end{lstlisting}
\end{summary}

\subsection*{Quick Experiments (Try It)}

\begin{tryit}{Order of traversals}
Build the tree: \texttt{A} at root; left child \texttt{B} with children \texttt{D} and \texttt{E};
right child \texttt{C} with right child \texttt{F}.
Print the preorder, inorder, and postorder sequences. Verify manually.
\end{tryit}

\begin{tryit}{Height of a skewed tree}
Build a right-chain: root$\to$1$\to$2$\to\cdots\to$9. Compute the height.
What is the height formula for $n$ nodes in a right-skewed tree?
Compare with a perfect binary tree of $n$ nodes.
\end{tryit}

\section*{In-Lab Exercises (target: 2 hours)}

\begin{labexercise}{Build and Traverse a Binary Tree \hfill {\small \itshape (45 min)}}
Implement \texttt{BinaryTree<E>} with:
\begin{itemize}[leftmargin=*,nosep]
  \item Methods: \texttt{root()}, \texttt{left(Node)}, \texttt{right(Node)}, \texttt{parent(Node)}.
  \item \texttt{addRoot(E e)}, \texttt{addLeft(Node p, E e)}, \texttt{addRight(Node p, E e)}.
  \item \texttt{int size()}, \texttt{boolean isEmpty()}.
  \item Traversals: \texttt{preorderPrint()}, \texttt{inorderPrint()}, \texttt{postorderPrint()}.
\end{itemize}
Construct a 7-node tree and verify all three traversal outputs.
\end{labexercise}

\begin{labexercise}{Height, Depth, and Leaf Count \hfill {\small \itshape (30 min)}}
Add:
\begin{itemize}[leftmargin=*,nosep]
  \item \texttt{int height()} --- height of the whole tree.
  \item \texttt{int depth(Node<E> v)} --- depth of node \texttt{v}.
  \item \texttt{int countLeaves()} --- nodes with no children.
  \item \texttt{int countNodes()} --- total nodes.
\end{itemize}
Test on at least two trees (balanced vs.\ skewed).
\end{labexercise}

\begin{labexercise}{Expression Tree Evaluator \hfill {\small \itshape (25 min)}}
Represent $(3 + 4) \times (8 - 5)$ as a binary tree (leaves = integers, internal = operators).
Write recursive \texttt{int evaluate(Node<String> v)}.
Print the inorder traversal (infix expression) and the evaluated result.
\end{labexercise}

\begin{labexercise}{Level-Order Traversal \hfill {\small \itshape (20 min)}}
Implement \texttt{void levelOrder()} using a Queue: enqueue root; while non-empty, dequeue a node,
visit it, enqueue non-null children. Print the output for your tree from Exercise 1.
\end{labexercise}

\begin{aicollab}
\textbf{Suggested prompts:}
\begin{itemize}[leftmargin=*,nosep]
  \item ``What is the difference between preorder, inorder, and postorder? Give a use case for each.''
  \item ``Explain why inorder traversal of a BST yields sorted order.''
\end{itemize}
\medskip\textbf{Verify:} Ask the AI for the three traversal sequences of a specific 7-node tree.
Build that tree in code and compare outputs.
\end{aicollab}

\takehomebanner

\begin{trace}[Traversal sequences]
Write the preorder, inorder, and postorder for the tree:
\begin{center}
\begin{tabular}{c}
\texttt{+} \\ \texttt{/ \textbackslash} \\ \texttt{*\quad 7} \\ \texttt{/ \textbackslash} \\ \texttt{3\quad 4}
\end{tabular}
\end{center}
(Root \texttt{+}; left subtree \texttt{*} with children \texttt{3} and \texttt{4}; right child \texttt{7}.)
\end{trace}

\begin{writecode}[Mirror a tree]
Write recursive \texttt{void mirror(Node<E> v)} that swaps left and right children of every node.
After \texttt{mirror(root)}, inorder should be the reverse of the original.
\end{writecode}

\begin{drawex}{Tree from traversals}
Reconstruct the unique binary tree with preorder $A, B, D, E, C, F$ and inorder $D, B, E, A, F, C$.
Draw the resulting tree and label every edge.
\end{drawex}

\vspace{4pt}
\noindent\textbf{Submission.} Save files as \texttt{lab8\_ex*.java} and submit on the course site.

\end{document}
